\documentclass{beamer}

\usetheme{Madrid}

\usepackage[utf8]{inputenc}
\usepackage[T5]{fontenc}
\usepackage[vietnamese]{babel}

\usepackage{amsmath}
\usepackage{tikz}
\usepackage{listings}
\usepackage{xcolor}
\usepackage{tcolorbox}
\usepackage{fontawesome5}

\definecolor{main1}{RGB}{102,126,234}
\definecolor{main2}{RGB}{118,75,162}

\setbeamercolor{frametitle}{fg=white,bg=main1}
\setbeamercolor{title}{fg=white,bg=main2}

\setbeamertemplate{background canvas}{
\begin{tikzpicture}[remember picture,overlay]
\path[left color=main1!20,right color=main2!20]
(current page.south west) rectangle (current page.north east);
\end{tikzpicture}
}

\tcbset{colback=white,colframe=main1,arc=4mm,boxrule=1pt,drop shadow}

\lstset{language=C++,basicstyle=\ttfamily\tiny,keywordstyle=\color{main2},commentstyle=\color{green!50!black},breaklines=true}

\title{Đường đi ngắn nhất trên đồ thị}
\author{Phạm Gia Kiệt}
\date{\today}

\begin{document}

\frame{\titlepage}

%================= INTRO =================
\begin{frame}{Bài toán đường đi ngắn nhất}

\begin{tcolorbox}
\textbf{\faRoute Khái niệm}
\begin{itemize}
\item Tìm đường đi từ đỉnh $u$ đến $v$ sao cho tổng trọng số nhỏ nhất
\item Có thể có nhiều đường đi khác nhau giữa hai đỉnh
\item Không phải đường trực tiếp luôn là tối ưu
\end{itemize}
\end{tcolorbox}

\vspace{0.2cm}

\begin{tcolorbox}
\textbf{\faLightbulb Ví dụ thực tế}
\begin{itemize}
\item Google Maps tìm đường đi nhanh nhất
\item Đi học: chọn đường ít tắc nhất
\item Mạng máy tính: truyền dữ liệu tối ưu
\end{itemize}
\end{tcolorbox}

\end{frame}

%================= TYPE =================
\begin{frame}{Phân loại bài toán}

\begin{tcolorbox}
\textbf{\faLayerGroup Các dạng phổ biến}
\begin{itemize}
\item Từ 1 đỉnh đến tất cả → Dijkstra
\item Giữa mọi cặp đỉnh → Floyd
\item Không trọng số → BFS
\end{itemize}
\end{tcolorbox}

\vspace{0.2cm}

\begin{tcolorbox}
\textbf{\faCheckCircle Ghi nhớ nhanh}
\begin{itemize}
\item Dijkstra dùng nhiều nhất trong thi
\item Floyd dùng khi $n$ nhỏ
\item BFS là trường hợp đặc biệt
\end{itemize}
\end{tcolorbox}

\end{frame}

%================= DIJKSTRA =================
\begin{frame}{Dijkstra - Ý tưởng}

\begin{tcolorbox}
\textbf{\faBolt Nguyên lý}
\begin{itemize}
\item Xuất phát từ đỉnh nguồn
\item Luôn chọn đỉnh có khoảng cách nhỏ nhất
\item Sau khi chọn thì giá trị đó là tối ưu
\end{itemize}
\end{tcolorbox}

\vspace{0.2cm}

\begin{tcolorbox}
\textbf{\faExclamationTriangle Lưu ý}
\begin{itemize}
\item Không dùng khi có cạnh âm
\item Nếu có cạnh âm → dùng Bellman-Ford
\end{itemize}
\end{tcolorbox}

\end{frame}

\begin{frame}{Dijkstra - Các bước}

\begin{tcolorbox}
\begin{enumerate}
\item Khởi tạo $d[s]=0$
\item Chọn đỉnh có khoảng cách nhỏ nhất
\item Cập nhật các đỉnh kề
\item Lặp lại đến khi hết
\end{enumerate}
\end{tcolorbox}

\vspace{0.2cm}

\begin{tcolorbox}
\textbf{\faClock Độ phức tạp}
\begin{itemize}
\item $O(E \log V)$
\item Chạy rất nhanh với đồ thị lớn
\end{itemize}
\end{tcolorbox}

\end{frame}

\begin{frame}[fragile]{Code Dijkstra}
\begin{tcolorbox}
\begin{lstlisting}
#define ll long long
const ll INF=1e18;
vector<pair<ll,ll>> adj[100005];
ll d[100005];

void f(ll s){
    for(ll i=1;i<=n;i++) d[i]=INF;
    priority_queue<pair<ll,ll>,vector<pair<ll,ll>>,greater<>> q;
    d[s]=0; q.push({0,s});
    while(q.size()){
        auto [du,u]=q.top(); q.pop();
        if(du!=d[u]) continue;
        for(auto [v,w]:adj)
            if(d[v]>du+w)
            d[v]=du+w,q.push({d[v],v});
    }
}
\end{lstlisting}
\end{tcolorbox}
\end{frame}

%================= FLOYD =================
\begin{frame}{Floyd - Ý tưởng}

\begin{tcolorbox}
\textbf{\faProjectDiagram Nguyên lý}
\begin{itemize}
\item Thử đi qua từng đỉnh trung gian k
\item Nếu đi qua k ngắn hơn → cập nhật
\item Áp dụng cho mọi cặp đỉnh
\end{itemize}
\end{tcolorbox}

\vspace{0.2cm}

\begin{tcolorbox}
\textbf{\faCalculator Công thức}
[
d[i][j] = \min(d[i][j], d[i][k] + d[k][j])
]
\end{tcolorbox}

\end{frame}

\begin{frame}[fragile]{Code Floyd}
\begin{tcolorbox}
\begin{lstlisting}
const ll INF=1e18;
ll d[505][505];

for(ll k=1;k<=n;k++)
    for(ll i=1;i<=n;i++)
        for(ll j=1;j<=n;j++)
            d[i][j]=min(d[i][j],d[i][k]+d[k][j]);
\end{lstlisting}
\end{tcolorbox}

\vspace{0.2cm}

\begin{tcolorbox}
\textbf{\faExclamationTriangle Lưu ý}
\begin{itemize}
\item Vòng k phải nằm ngoài cùng
\item Không dùng khi có chu trình âm
\end{itemize}
\end{tcolorbox}

\end{frame}

%================= COMPARE =================
\begin{frame}{So sánh}

\begin{tcolorbox}
\begin{tabular}{|c|c|c|}
\hline
Thuật toán & Độ phức tạp & Dùng cho \
\hline
Dijkstra & $O(E\log V)$ & 1 nguồn \
Floyd & $O(n^3)$ & Mọi cặp \
\hline
\end{tabular}
\end{tcolorbox}

\vspace{0.2cm}

\begin{tcolorbox}
\textbf{\faLightbulb Kết luận nhanh}
\begin{itemize}
\item Dijkstra: nhanh, dùng nhiều nhất
\item Floyd: đơn giản nhưng chậm
\item Chọn đúng thuật toán là quan trọng
\end{itemize}
\end{tcolorbox}

\end{frame}

\begin{frame}
\centering
\Huge{Cảm ơn!}
\end{frame}

\end{document}
